\chapter{Metodologi Penelitian}

\section{Alur Penelitian}
Penelitian ini dijalankan secara sistematis mengikuti tahapan yang terstruktur, mulai dari identifikasi masalah, studi literatur, pengumpulan data, konstruksi model, hingga tahap evaluasi dan penarikan kesimpulan. Secara visual, alur ini direpresentasikan dalam sebuah \textit{flowchart} penelitian.

\section{Pengumpulan Data \& Pra-pemrosesan}
Data yang digunakan dalam penelitian ini adalah data sekunder berupa harga historis aset yang tergabung dalam indeks LQ45.
\begin{itemize}
    \item \textbf{Sumber Data:} Akuisisi data dilakukan melalui API Yahoo Finance.
    \vspace{0.5em}
    \begin{itemize}
        \item[] \textbf{Algoritma Akuisisi Data:}
        \item[-] \textbf{Input:} Daftar \textit{ticker} saham ($T$), Tanggal mulai ($t_{start}$), Tanggal akhir ($t_{end}$).
        \item[-] \textbf{Proses:} 
            \begin{enumerate}[label=\arabic*.]
                \item Unduh harga harian ($P$).
                \item Hapus nilai kosong (\textit{drop NaN}).
                \item Urutkan berdasarkan waktu.
            \end{enumerate}
        \item[-] \textbf{Output:} Matriks \textit{time-series} harga ($P$).
    \end{itemize}
    \vspace{0.5em}
    \item \textbf{Digitalisasi Data Finansial:} Transformasi \textit{log return} kontinu menjadi \textit{state} biner menggunakan teknik diskretisasi sebagai jembatan menuju sistem qubit.
    \item \textbf{Laplace Smoothing:} Diterapkan untuk menstabilkan perhitungan probabilitas pada estimasi entropi dan korelasi.
\end{itemize}

\section{Konstruksi Model Permainan (Mapping Pipeline)}
Tahap ini menjelaskan bagaimana data finansial dipetakan ke dalam model \textit{Quantum Game Theory}.
\begin{itemize}
    \item \textbf{Dinamika Parameter Risiko Endogen:} Penentuan parameter risiko ($\lambda$) yang adaptif.
    \vspace{0.5em}
    \begin{itemize}
        \item[] \textbf{Algoritma Kalkulasi $\lambda$:}
        \item[-] \textbf{Input:} \textit{Log return} historis ($R$).
        \item[-] \textbf{Proses:}
            \begin{enumerate}[label=\arabic*.]
                \item Hitung rata-rata \textit{return} tahunan absolut ($\mu_{avg}$) dan volatilitas tahunan ($\sigma_{avg}$).
                \item Evaluasi rasio performa: $Z = \frac{\mu_{avg}}{\sigma_{avg}}$.
                \item Terapkan fungsi \textit{sigmoid}: $\lambda = \frac{1}{1 + e^Z}$.
            \end{enumerate}
        \item[-] \textbf{Output:} Parameter risiko adaptif $\lambda \in (0, 1)$.
    \end{itemize}
    \vspace{0.5em}

    \item \textbf{Matriks Payoff:} Perhitungan nilai utilitas menggunakan kombinasi \textit{return} dan penalti volatilitas berbobot $\lambda$.
    \vspace{0.5em}
    \begin{itemize}
        \item[] \textbf{Algoritma Kalkulasi Matriks Payoff Ekspektasi:}
        \item[-] \textbf{Input:} Rangkaian \textit{state} aksi biner ($s_A, s_B \in \{0, 1\}$), \textit{return} ($R_A, R_B$), dan $\lambda$.
        \item[-] \textbf{Proses:}
            \begin{enumerate}[label=\arabic*.]
                \item Inisialisasi matriks \textit{payoff} ($p_A, p_B$) berukuran $2 \times 2$ dengan nilai 0.
                \item Untuk setiap observasi waktu $t$:
                \begin{itemize}
                    \item Identifikasi profil kejadian $(i, j) = (s_{A, t}, s_{B, t})$.
                    \item Hitung utilitas sesaat: $u_{A, t} = (1-\lambda)R_{A, t} - \lambda |R_{A, t}|$.
                    \item Akumulasi ke elemen matriks: $p_A[i, j] \mathrel{+}= u_{A, t}$ dan tambah penghitung frekuensi.
                \end{itemize}
                \item Normalisasi setiap sel $p_A[i, j]$ dengan total frekuensi kejadian $(i, j)$.
            \end{enumerate}
        \item[-] \textbf{Output:} Matriks \textit{payoff} pemain ($p_A, p_B$).
    \end{itemize}
    \vspace{0.5em}

    \item \textbf{Dinamika Stackelberg:} Identifikasi peran \textit{Leader} dan \textit{Follower} berdasarkan kapitalisasi pasar untuk memodelkan asimetri informasi.
    
    \item \textbf{Ekstraksi Parameter Hamiltonian:} Menentukan nilai bias laten ($h_i$) dari margin \textit{payoff} dan nilai kopling ($J_{ij}$) berdasarkan ukuran jarak entropik (\textit{Quantum Mutual Information}).
    \vspace{0.5em}
    \begin{itemize}
        \item[] \textbf{Algoritma Kalkulasi Hamiltonian $h_i$ (Bias):}
        \item[-] \textbf{Input:} Matriks \textit{payoff} seluruh pasangan aset, jumlah total aset ($N$).
        \item[-] \textbf{Proses:}
            \begin{enumerate}[label=\arabic*.]
                \item Untuk setiap iterasi agregasi pada aset ke-$i$:
                \item Hitung selisih \textit{payoff} marginal dari probabilitas koordinasi matriks $p$: $(p[1,0] + p[1,1]) - (p[0,0] + p[0,1])$.
                \item Rata-ratakan selisih dari seluruh pasangan yang melibatkan aset $i$.
            \end{enumerate}
        \item[-] \textbf{Output:} Vektor parameter bias sentral $h_i$.
    \end{itemize}
    \vspace{0.5em}
    \begin{itemize}
        \item[] \textbf{Algoritma Kalkulasi Interaksi QMI ($J_{ij}$):}
        \item[-] \textbf{Input:} \textit{State} biner aset $A$ dan $B$ ($s_A, s_B$).
        \item[-] \textbf{Proses:}
            \begin{enumerate}[label=\arabic*.]
                \item Hitung probabilitas gabungan \textit{Laplace-smoothed} untuk setiap kuadran keadaan $\rho_{AB} = \text{diag}(\dots)$.
                \item Lakukan reduksi untuk mendapatkan \textsl{reduced density matrix} marjinal $\rho_A = \text{tr}_B(\rho_{AB})$ dan $\rho_B = \text{tr}_A(\rho_{AB})$.
                \item Hitung entropi Von Neumann: $S(\rho) = -\sum \lambda_k \ln(\lambda_k)$ untuk \textit{eigenvalue} $\lambda_k > 0$.
                \item Ekstrak QMI: $J_{ij} = S(\rho_A) + S(\rho_B) - S(\rho_{AB})$.
            \end{enumerate}
        \item[-] \textbf{Output:} Nilai kopling keterikatan $J_{ij}$.
    \end{itemize}
    \vspace{0.5em}

    \item \textbf{Klasifikasi Rezim:} Identifikasi apakah interaksi aset cenderung bersifat kooperatif (\textit{Coordination}) atau kompetitif (\textit{Anti-Coordination}).
\end{itemize}

\section{Implementasi Kuantum \& Optimasi}
Proses optimasi dilakukan menggunakan algoritma VQE pada simulator kuantum.
\begin{itemize}
    \item \textbf{Formulasi Hamiltonian Pembatas:} Agregasi operator observasi kuantum menjadi bentuk energi total dengan batasan kardinalitas.
    \vspace{0.5em}
    \begin{itemize}
        \item[] \textbf{Algoritma Penempatan Hamiltonian:}
        \item[-] \textbf{Input:} Vektor bias $h$, matriks interaksi $J$, batas alokasi $K$, faktor penalti $A$.
        \item[-] \textbf{Proses:}
            \begin{enumerate}[label=\arabic*.]
                \item Tempatkan operator lokal Pauli-$Z_i$ berbobot skalar $h_i$ untuk setiap basis aset.
                \item Tempatkan pasangan operator ikatan berbobot $J_{ij}Z_i Z_j$.
                \item Injeksi konstanta pembatas global $A \times I$.
                \item Aplikasikan penalti pada setiap elemen sekuensial $\frac{A}{2} Z_i Z_j$.
            \end{enumerate}
        \item[-] \textbf{Output:} Hamiltonian komputasional $H$.
    \end{itemize}
    \vspace{0.5em}

    \item \textbf{Konfigurasi dan Eksekusi VQE dengan Optimasi SPSA:} Transformasi ising ke sirkuit parametrik kuantum hingga mendapatkan konfigurasi optimal.
    \vspace{0.5em}
    \begin{itemize}
        \item[] \textbf{Algoritma VQE-SPSA:}
        \item[-] \textbf{Input:} Operator Hamiltonian $H$, jumlah \textit{qubits} (aset), batas kedalaman \textit{ansatz} (\textit{depth}=1).
        \item[-] \textbf{Proses:}
            \begin{enumerate}[label=\arabic*.]
                \item Deklarasikan arsitektur variatif (\textit{hardware-efficient ansatz}): Rotasi $RY$ dan $RZ$ diikuti cincin \textit{entanglement} $CNOT$.
                \item Inisialisasi sudut rotasi $\vec{\theta}$ secara acak $\mathcal{U}(0, 2\pi)$.
                \item Untuk siklus evaluasi konvergensi $k$:
                \begin{itemize}
                    \item Bangkitkan arah perturbasi stokastik biner $\vec{\Delta}_{k} \in \{-1, +1\}^m$.
                    \item Evaluasi gradien aproksimasi simetris: 
                    $g_k = \frac{\langle H(\vec{\theta} + c_k\vec{\Delta}_{k}) \rangle - \langle H(\vec{\theta} - c_k\vec{\Delta}_{k}) \rangle}{2 c_k \vec{\Delta}_{k}}$
                    \item Perbarui sudut rotasi: $\vec{\theta} \mathrel{-}= a_k g_k$.
                \end{itemize}
                \item Ekstraksi probabilitas densitas pada bentuk gelombang (\textit{probability distribution}) dengan jumlah aset dipilih $= K$.
            \end{enumerate}
        \item[-] \textbf{Output:} Subset indeks aset bernilai \textit{state} `1' berprobabilitas tertinggi.
    \end{itemize}
\end{itemize}

\section{Evaluasi dan Analisis Kinerja}
Hasil optimasi dievaluasi untuk mengukur efektivitas model.
\begin{itemize}
    \item \textbf{Identifikasi State Optimal:} Ekstraksi \textit{bitstring} dengan probabilitas tertinggi sebagai representasi alokasi aset terbaik.
    \item \textbf{Backtesting \& Slidng Window:} Analisis historis untuk mengamati perubahan alokasi aset terhadap peristiwa makroekonomi (\textit{regime shift}).
    \item \textbf{Validasi Metrik:} Perbandingan kinerja antara strategi berbasis kuantum dengan strategi klasik menggunakan metrik \textit{Sharpe Ratio} dan \textit{Trend Ratio}.
\end{itemize}

\newpage
\thispagestyle{empty}
\vspace*{\fill}
    \begin{center}
	Halaman ini sengaja dikosongkan
    \end{center}
\vspace*{\fill}
\pagebreak
